\documentclass[11pt]{article}
\usepackage[T1]{fontenc}
\usepackage[a4paper,margin=1in]{geometry}
\usepackage{amsmath,amssymb,amsthm}
\usepackage[hidelinks]{hyperref}
\usepackage{microtype}

\newtheorem*{question}{Source Problem 3.12}
\newtheorem*{answer}{Later result}

\setlength{\parindent}{0pt}
\setlength{\parskip}{0.55em}

\title{Literature Status: Nonlinear 2-Local\\
$*$-Homomorphisms of $C^*$-Algebras}
\author{Run \texttt{fa\_banach\_001} --- lane 0}
\date{9 August 2026}

\begin{document}
\maketitle

\begin{center}
\fbox{\parbox{0.9\textwidth}{
\textbf{Status: literature already answered (partial subcase).}
The $*$-homomorphism half of Problem~3.12 in arXiv:1312.6850 was
explicitly answered by Cabello--Peralta, arXiv:1510.00915,
Corollary~2.8.  The plain non-$*$ homomorphism half is not covered.
}}
\end{center}

\section{Original question}

Antonio M. Peralta's \emph{A note on 2-local representations of
$C^*$-algebras}, arXiv:1312.6850, later published in
\emph{Operators and Matrices} \textbf{9}(2) (2015), 343--358,
poses the following as Problem~3.12 on arXiv PDF page~12~\cite{source}.

\begin{question}
Is every (not necessarily linear) 2-local ($*$-)homomorphism between
$C^*$-algebras a ($*$-)homomorphism? Equivalently, is the linearity
hypothesis in Theorem~3.9 superfluous?
\end{question}

The notation simultaneously asks two questions: one for bounded algebra
homomorphisms used as the two-point interpolants, and one for
$*$-homomorphisms.  Source Theorem~3.9 already proves the conclusion once
the 2-local map itself is linear (and bounded in the non-$*$ case).

\section{Separate later answer to the \texorpdfstring{$*$}{star}-half}

Juan Carlos Cabello and Antonio M. Peralta's separate later paper,
\emph{Weak-2-local symmetric maps on $C^*$-algebras}, arXiv:1510.00915,
published in \emph{Linear Algebra and its Applications} \textbf{494}
(2016), 32--43, proves automatic linearity for a larger symmetric class
\cite{answerpaper}.  It then specializes this result on arXiv PDF page~8:

\begin{answer}
\leavevmode\par\noindent
Every 2-local $*$-homomorphism between $C^*$-algebras is linear and is a
$*$-homomorphism.
\end{answer}

The citation trail is exact.  Corollary~2.7 gives linearity.  The paragraph
immediately before Corollary~2.8 combines that corollary with Theorem~3.9 of
the source paper, cited there as reference~[27], and says this yields the
positive general result.  Thus the later authors explicitly use the source
theorem to remove its linearity hypothesis for 2-local $*$-homomorphisms.

\section{Scope and search status}

This settles the $*$-preserving half of Problem~3.12 for arbitrary source and
target $C^*$-algebras.  It does not settle the plain homomorphism half:
the automatic-linearity mechanism in the supporting paper applies to
symmetric maps, and an arbitrary homomorphism need not preserve involution.

A bounded search through 9 August 2026 checked the run indexes, the exact
wording of Problem~3.12, arXiv searches for nonlinear 2-local homomorphisms
between $C^*$-algebras, and the source paper's citation trail.  The decisive
match was arXiv:1510.00915.  Searches also recovered the earlier
von-Neumann/compact-$C^*$ result arXiv:1404.7597, but no general result for
the plain non-$*$ half.  No new mathematical result is claimed here.

\begin{thebibliography}{9}
\bibitem{source}
A. M. Peralta,
\emph{A note on 2-local representations of $C^*$-algebras},
Operators and Matrices \textbf{9}(2) (2015), 343--358;
arXiv:1312.6850.  doi:10.7153/oam-09-20.

\bibitem{answerpaper}
J. C. Cabello and A. M. Peralta,
\emph{Weak-2-local symmetric maps on $C^*$-algebras},
Linear Algebra Appl. \textbf{494} (2016), 32--43;
arXiv:1510.00915.  doi:10.1016/j.laa.2015.12.024.
\end{thebibliography}

\end{document}
